\section{文书、题记与遗址：南北朝社会的可见与不可见}
\label{sec:northern-southern-documents-material}

南北朝的帝纪能够排出受禅、战争与诏令的顺序，却很少完整保留一座城市如何筹料、一户人家如何迁徙或一处寺院如何经营。四四〇年前后沮渠无讳率北凉遗众西趋高昌，四四二年又在绿洲城市吸纳河西随从；《魏书》〈高昌传〉与吐鲁番文书、墓葬考古使我们能看见河西流亡集团、粮食和行政资源的某些线索。使节往来和人群迁移可相互校读，族称、疆界和城市内部的控制程度却不能由中原王朝的传记或一份文书单独固定。

\subsection{石窟不是帝国的缩影}

四四六年北魏禁佛、四五二年恢复佛教，制度转折可见于《魏书》〈释老志〉；云冈、龙门的题记与考古为政策的物质环境增加另一种证词。四五〇年代昙曜主持云冈早期大像，石工、供养人和物资被组织进平城周边的营建；四六九年延续的洞窟与题记又显示供养关系不只来自一个层级。题记可校验部分姓名、日期和功德表达，洞窟与造像主体的年代却未必完全一致，更不能由王权赞助或一处造像推断全境信仰。

禁佛的地方边界同样须由这些材料来限制。北魏诏令确实使僧尼、寺产和造像受到冲击，但军镇与州郡的执行程度不一；文本与石窟材料相合，只足以说明毁禁和复兴需要分地区考察。北周五七四年禁佛时，寺产、僧尼与宗教器物被清查、再分配；北齐仍有寺院与造像赞助，灭齐后又出现不同幅度的恢复。墓志、造像记和遗址能使这些差异不被诏令抹平，仍不能替我们计算寺院资产或断言某一地区的全体信众如何回应。

\subsection{新都的物质条件与怀旧的文字}

四九四年迁都后，洛阳城郭、宫城、道路、坊市吸纳官营资源、工匠和迁入人口；五一六年永宁寺及高塔的营建又需要木材、运输和劳役。都城考古、题记与《洛阳伽蓝记》能补足本纪对工匠和供给的省略，却各有不同的时间层次和写作目的。特别是杨衒之约五四七年撰成《洛阳伽蓝记》，是在旧都衰败、政权已迁移之后追述繁盛与毁乱；它是理解洛阳寺院、城市记忆和佛教网络的重要证词，不能无条件当作五一六年工地的现场账簿。

同一限制也适用于南方。侯景之乱破坏建康寺院、市场、漕运和户籍，陈朝接收梁末州郡文书、军人和侨民后才逐步重整周边秩序。正史、建康遗址和墓志能提示行政与城市社会的断裂，却很难连续追踪普通住民在围城、迁徙与恢复中的经历。将都城的可见遗存直接等同于整个江南，反而会重演史料本身偏向精英与中枢的局限。

\subsection{证据的交叉，不是证据的相加}

《魏书》《北齐书》《周书》与南朝诸史提供制度和政权年序；吐鲁番文书、墓志、墓券、造像记、石窟寺院遗址和都城考古，分别提供地方行政、迁徙、功德、技术与空间的有限锚点。墓志的官衔、世系和道德称颂常带有家族自我呈现，题记也可能残缺、重刻或脱离原位；考古分期不能独自证明具体政治事件。把材料并读的目的，不是把每一种证词简单相加成完整社会，而是区分能够互证的日期、姓名与空间，同仍然无从量化的户口、信仰、劳役和动机。这样理解，南北朝的文书与遗址既扩展了正史的视野，也使我们更准确地看见它们不能照亮的部分。
